\section{Field Extensions}

We build up the theory of field extensions following Morandi~\cite{Morandi1996}.

The study of field theory is primarily the study of field extensions. The classical problems of ruler-and-compass constructions and the solvability of polynomial equations by radicals were settled by analyzing appropriate field extensions. The central idea is to associate a group---the \emph{Galois group}---to a field extension, converting field-theoretic questions into group-theoretic ones. Since the Galois group of a finite extension is finite, the rich structure theory of finite groups becomes available. In this section we build the foundational vocabulary: degrees, algebraic elements, minimal polynomials, and the key structural results about simple and algebraic extensions.

%------------------------------------------------------------------
\subsection{Degrees and Generated Subfields}
%------------------------------------------------------------------
Let $F\subseteq K$ are fields. 
Then $K$ is called a \emph{field extension} of $F$. 
We denote it by $K/F$. 
Here $F$ is called the \emph{base field}. 
We can think of $K$ as a vector space over $F$ where the addition of two vectors follows the addition operation of $K$ and the multiplication of any $\alpha\in K$ and $c\in F$ follows the multiplication operation of $K$, $c\cdot \alpha=c\alpha$. 
Then we denote the dimension of $K$ over $F$ by  $$[K\colon F]\coloneqq\dim_F(K)$$This dimension is called the \emph{degree of field extension} of $F$ or simply \emph{degree of} $[K\colon F]$. 
If $[K\colon F]<\infty$ we call it by a special name, \emph{finite extension}. 
Finite extensions of finite fields are of very important to us and this report we will work on finite fields and their finite extensions. 
Let $\alpha\in K$. 
Then we denote $F(\alpha)$ to be the field generated by $F$ and $\alpha$ following the multiplication and addition operations of $K$. 
We will define what it means to be the field generated by $F$ and $\alpha$ in \defref{def:generated-subfields}.
It is very easy to sea $F(\alpha)$ is also a field, and it contains $F$ but also contained in $K$. 

Let $F$ be any field and suppose $X$ be a formal variable. 
Then $F[X]$ denotes the polynomial ring over $F$. 
We define the rational function field $F(X)$ to be the \emph{fraction ring} of $F[X]$ which consists of all quotients $\nicefrac{f}{g}$ where $f,g\in F[X]$ with $g(X)\neq 0$.
Similarly, we can define the polynomial ring $F[X_1,\dots, X_n]$ over multiple formal variables and their fraction field $F(X_1,\dots, X_n)$.

Look at the following example. $\bbR\subseteq \bbC$ are both fields. But $\bbC$ is a $2$-dimensional vector space over $\bbR$ as we can take the basis $\{1,i\}$.
So $\bbC/\bbR$ is a finite extension of degree $2$. 
But $\bbQ/\bbR$ is an infinite field extension. 
On the other hand let $\alpha\in \bbR$. Then consider the field $\bbQ(\alpha)$. 
Now degree of  $[\bbQ(\alpha)\colon \bbQ]$ depends on what kind of element $\alpha$ is. If $\alpha=\sqrt{2}$ then $[\bbQ(\sqrt{2})\colon \bbQ]=2$ but if $\alpha=\pi$ then $[\bbQ(\pi)\colon \bbQ]=\infty$.

\begin{Definition}{Generated Subfields}{def:generated-subfields}
  Let $K/F$ be a field extension and $S\subseteq K$.
  \begin{enumerate}[label=(\roman*)]
    \item The \emph{ring generated by $F$ and $S$}, written $F[S]$, is the intersection of all subrings of $K$ containing $F$ and $S$.
    \item The \emph{field generated by $F$ and $S$}, written $F(S)$, is the intersection of all subfields of $K$ containing $F$ and $S$.
  \end{enumerate}
  If $S=\{a_1,\ldots,a_n\}$ is finite we write $F[a_1,\ldots,a_n]$ and $F(a_1,\ldots,a_n)$. 
  We call $F(\alpha_1,\dots, \alpha_n)$ \emph{finitely generated} field. 
\end{Definition}
From this definition it easily follows that for any $S\subseteq K$, $F(S)$ is the smallest subfield of $K$ containing $F$ and $\alpha$.

Let $K/F$ be a field extension. 
For any $\alpha\in K$ consider the map $\ev_{\alpha}:F[X]\to K$ defined by $f(X)\mapsto f(\alpha)$. 
Hence, $\ev_{\alpha}$ is a ring homomorphism.
\begin{observation}
  For any $\alpha\in K$, $\image(\ev_{\alpha})$ is a subring of $K$ and $\ev_{\alpha}$ is $F$-vector space homomorphism.
\end{observation}

Now we can use this map $\ev_{\alpha}$ to give descriptions of the elements of $F[\alpha]$ and $F(\alpha)$. 
\begin{lemma}{}{lem:simple-gen}
  Let $K$ be a field extension of $F$ and let $a\in K$. Then
  \[
    F[a] = \{f(a) : f(x)\in F[x]\}
    \quad\text{and}\quad
    F(a) = \left\{\frac{f(a)}{g(a)} : f,g\in F[x],\; g(a)\neq 0\right\}.
  \]
  Moreover, $F(a)$ is the quotient field of $F[a]$.
\end{lemma}

\begin{proof}
  The evaluation map $\mathrm{ev}_a$ is a ring homomorphism with image $\{f(a):f\in F[x]\}$/ 
  So this set is a subring of $K$. 
  Any subring of $K$ containing $F$ and $a$ must contain $f(a)$, for all $f\in F[X]$. 
  Therefore, $F[a]=\{f(a):f\in F[x]\}$. 
  The fraction field of $F[a]$ is $\{f(a)/g(a)\colon g(a)\neq 0\}$; it is contained in any subfield of $K$ containing $F[a]$, hence equals $F(a)$.
\end{proof}

We call this procedure \emph{adjoining $\alpha$} to $F$. 
Using similar arguments we can describe the ring $F[a_1,\dots, a_n]$ and the field $F(a_1,\dots, a_n)$ where $K/F$ is a field extension and $a_1,\dots, a_n\in K$. 
So we have the following theorem.
\begin{Theorem}{}{thm:mult-gen}
  Let $K$ be a field extension of $F$ and let $a_1,\ldots,a_n\in K$. Then
  \begin{align*}
    F[a_1,\ldots,a_n] &= \{f(a_1,\ldots,a_n) : f\in F[x_1,\ldots,x_n]\}, \\
    F(a_1,\ldots,a_n) &= \left\{\frac{f(a_1,\ldots,a_n)}{g(a_1,\ldots,a_n)} : f,g\in F[x_1,\ldots,x_n],\; g(a_1,\ldots,a_n)\neq 0\right\}.
  \end{align*}
  Moreover, $F(a_1,\ldots,a_n)$ is the quotient field of $F[a_1,\ldots,a_n]$.
\end{Theorem}

Now we can also define fields generated by adjoining arbitrary number of elements. 
In this case we can give the description of the new field in terms of all finite extension of fields created by adjoining finitely many elements. 

\begin{Theorem}{}{thm:fin-gen-sub}
  Let $K$ be a field extension of $F$ and let $S\subseteq K$ be any subset. 
  If $\alpha\in F(S)$, then $\alpha\in F(a_1,\ldots,a_n)$ for some $a_1,\ldots,a_n\in S$. 
  Therefore,
  \[
    F(S) = \bigcup\,\{F(a_1,\ldots,a_n) \colon a_1,\ldots,a_n\in S\},
  \]
  where the union is over all finite subsets of $S$.
\end{Theorem}

\begin{proof}
  Let $U\coloneqq \cup\,\{F(a_1,\ldots,a_n) \colon a_1,\ldots,a_n\in S\}$. 
  Each $F(a_1,\ldots,a_n)$ contains $F$ and lies in $F(S)$ where $a_1,\dots, a_n\in S$. 
  So the union $U$ satisfies $U\subseteq F(S)$. 
  The set $U$ contains $F$ and $S$. 
  If $\alpha\in F(a_1,\ldots,a_n)$ and $\beta\in F(b_1,\ldots,b_m)$, then $\alpha\pm\beta,\,\alpha\beta,\,\alpha/\beta$ (if $\beta\neq 0$) are all in $F(a_1,\ldots,a_n,b_1,\ldots,b_m)\subseteq U$. 
  Therefore, $U$ is a field. Being the smallest subfield containing $F$ and $S$, it equals $F(S)$.
\end{proof}

%------------------------------------------------------------------
\subsection{Algebraic Elements and Minimal Polynomials}
%------------------------------------------------------------------

\begin{Definition}{Algebraic and Transcendental Elements}{algebraic-elem}
  Let $K$ be a field extension of $F$. 
  An element $\alpha\in K$ is \emph{algebraic over $F$} if there is a nonzero polynomial $f(X)\in F[X]$ with $f(\alpha)=0$. 
  If $\alpha$ is not algebraic over $F$, then $\alpha$ is called \emph{transcendental over $F$}. 
  If every element of $K$ is algebraic over $F$, then $K$ is said to be \emph{algebraic over $F$}, and $K/F$ is called an \emph{algebraic extension}.
\end{Definition}

Hermite (1873) proved that $e$ is transcendental over $\mathbb{Q}$ and Lindemann (1882) proved the same for $\pi$. 
On the other hand, $\sqrt[n]{r}\in\mathbb{R}$ is algebraic over $\mathbb{Q}$ for any $r\in\mathbb{Q}$ (root of $x^n-r$), as is $i=\sqrt{-1}$ (root of $x^2+1$).

\begin{Definition}{Minimal Polynomial}{min-poly}
  If $\alpha$ is algebraic over a field $F$, the \emph{minimal polynomial} of $\alpha$ over $F$ is the monic polynomial $p(X)$ of least degree in $F[X]$ for which $p(\alpha)=0$; it is denoted $\min(F,\alpha)$. 
  Equivalently, $\min(F,\alpha)$ is the monic generator of the kernel of the evaluation homomorphism $\mathrm{ev}_\alpha$.
\end{Definition}

The minimal polynomial depends on the base field. 
For instance, $\min(\mathbb{Q},i)=x^2+1$ while $\min(\mathbb{C},i)=x-i$, since $i\in\mathbb{C}$ is already in the base field.

The minimal polynomial of an element and the degree of a field extension
are two of the most basic tools we shall use. 
The following proposition gives a relation between these objects.

\begin{Theorem}{}{thm:simple-ext}
  Let $K/F$ be a field extension and let $\alpha\in K$ be algebraic over $F$. 
  If $p(X)=\min(F,\alpha)$, then:
  \begin{enumerate}[label=\textup{(\roman*)}]
    \item $p(X)$ is irreducible over $F$.
    \item If $g(X)\in F[X]$, then $g(\alpha)=0$ if and only if $p(X)\mid g(X)$.
    \item If $n=\deg(\min(F,\alpha))$, then the elements $1,\alpha,\ldots,\alpha^{n-1}$ form a basis for $F(\alpha)$ over $F$. 
    So $[F(\alpha)\colon F]=\deg(\min(F,\alpha))<\infty$. Moreover, $F(\alpha)=F[\alpha]$.
    \item $F(\alpha)\cong \quotient{F[X]}{\la p(X)\ra}$.
  \end{enumerate}
\end{Theorem}
\begin{proof}
  The \emph{first isomorphism theorem} gives $F[\alpha]\cong \quotient{F[X]}{\la p(X)\ra}$. 
  Since $F[\alpha]\subseteq K$ is an integral domain, $\langle p(X)\rangle$ is a prime ideal. 
  Now in a \pid every nonzero prime is maximal, so $\quotient{F[X]}{\la p(X)\ra}$ is a field and $F[\alpha]=F(\alpha)$. 
  This gives (i) and (iv).  
  Part (ii) follows from $\ker(\ev_\alpha)=\langle p(X)\rangle$.

  For (iii): any $b\in F(\alpha)=F[\alpha]$ satisfies $b=g(\alpha)$. 
  Dividing gives $g=qp+r$ with $\deg r<n$ and $q,r\in F[X]$. 
  So $b=r(\alpha)$ lies in the span of $\{1,\alpha,\ldots,\alpha^{n-1}\}$. 
  Therefore, the set $\{1,\alpha,\dots, \alpha^{n-1}\}$ spans $F(\alpha)$ as a vector space. 
  If $a_0+a_1\cdot \alpha+\cdots +a_{n-1}\cdot \alpha^{n-1}=0$ where $a_0,\dots, a_{n-1}\in F$ then the polynomial $f(X)=\sum_{i=0}^{n-1}a_iX^i$ is in $\ker(\ev_\alpha)$. 
  Hence, either $f\equiv 0$ or $p\mid f$ by part (ii). 
  But $\deg(f)<n$. 
  So $f\equiv 0$. 
  Thus, $1,\alpha,\dots, \alpha^{n-1}$ forms a basis of $F(\alpha)$ over $F$. 
\end{proof}

The part (iii) of above theorem gives us a basis of $F(\alpha)$ over $F$ using only $\alpha$. 
So we get the following observation:
\begin{observation}
  If $K/F$ is a field extension and $\alpha\in K$ is algebraic over $F$ where $p=\min(F,\alpha)$ then $[F(\alpha)\colon F]=\deg(p)$. 
\end{observation}

Notice that if $\alpha,\beta\in K$ are both algebraic over $F$ and $\min(F,\alpha)=\min(F,\beta)=p$ i.e., their minimal polynomials are same then by the above theorem $$F(\alpha)\cong\quotient{F[X]}{\la p\ra}\cong F(\beta)$$
Therefore we have the observation:
\begin{observation}
  If $K/F$ is a field extension and $\alpha,\beta\in K$ are both algebraic over $F$, and they have the same minimal polynomials then $F(\alpha)\cong F(\beta)$ with the isomorphism $\varphi:\alpha\mapsto \beta$. 
\end{observation}
%------------------------------------------------------------------
\subsection{Finite and Algebraic Extensions}
%------------------------------------------------------------------

\begin{lemma}{}{thm:finite-alg}
  If $[K:F]<\infty$, then $K$ is algebraic and finitely generated over $F$.
\end{lemma}

\begin{proof}
  Let $\{\alpha_1,\ldots,\alpha_n\}$ be an $F$-basis of $K$. 
  Then for any element $\alpha\in K$, there exists $a_1,\dots, a_n$ such that $\alpha=a_1\cdot \alpha_1+\cdots +a_n\cdot \alpha_n$. 
  So certainly we have $K=F(\alpha_1,\dots, \alpha_n)$ and hence $K$ is finitely generated over $F$. 
  Since $[K\colon F]=n$, for any $\alpha\in K$, the $n+1$ elements $1,\alpha,\ldots,\alpha^n$ span only an $n$-dimensional space. 
  So they are $F$-linearly dependent. 
  Thus, there exists $\beta_0,\dots, \beta_n\in F$ such that $\sum_{i=0}^n\beta_i\cdot \alpha^i=0$. So we get a polynomial $f(X)=\sum_{i=0}^n \beta_i\cdot X^i$ and $f(\alpha)=0$. 
  Therefore, $\alpha$ is algebraic over $F$. 
\end{proof}

Let $K/F$ is a finite extension and $K=F(\alpha_1,\dots, \alpha_n)$.  
Then we can break up the extension $K/F$ into collection of subextensions. 
Let $L_i=F(\alpha_1,\dots,\alpha_i)$. 
Set $L_0=F$. 
Then basically $L_{i+1}=L_i(\alpha_{i+1})$. 
And we have $$F=L_0\subseteq L_1\subseteq \cdots \subseteq L_{n-1}\subseteq L_n=K$$
We will show a theorem which relates the degree of the extension $K/F$ and each subextensions. 

\begin{Theorem}{}{thm:tower-law}
  Let $F\subseteq L\subseteq K$ be fields. Then $[K:F]=[K:L]\cdot[L:F]$.
\end{Theorem}

\begin{proof}
  Let $\{a_i\}_{i\in I}$ be an $F$-basis of $L$ and $\{b_j\}_{j\in J}$ an $L$-basis of $K$. 
  We show $\{a_i\cdot b_j\}_{I\times J}$ is an $F$-basis of $K$. 
  For any $\gm\in K$, write $\gm=\sum_j\alpha_jb_j$ with $\alpha_j\in L$. 
  Now each $\alpha_j=\sum_i\beta_{ij}a_i$ with $\beta_{ij}\in F$; then $\gm=\sum_{i,j}\beta_{ij}a_ib_j$. 
  Therefore, the set $\{a_i\cdot b_j\}_{I\times J}$ spans $K$ over $F$. 
  
  If $\sum_{i,j}\beta_{ij}a_ib_j=0$ for some non-trivial linear combinations where $\beta_{i,j}\in F$ for all $i\in I$, $j\in J$. 
  Now $$\sum_{i,j}\beta_{ij}a_ib_j=\sum_j\bigl(\sum_i\beta_{ij}a_i\bigr)b_j=0.$$ 
  Since $\{b_j\}_{j\in J}$ is a basis of $K$ over $L$, each $\sum_i\beta_{ij}a_i=0$. 
  And by $F$-independence of $\{a_i\}$, each $\beta_{ij}=0$. 
  So $\{a_i\cdot b_j\}_{I\times J}$ form a $F$-basis of $K$. 
  Therefore, $$[K:F]=|\{a_i\cdot b_j\}_{I\times J}|=|\{a_i\colon i\in I\}|\cdot |\{b_j\colon j\in J\}|=[K\colon L]\cdot [L\colon F]$$ 
  So we have the theorem. 
\end{proof}

\begin{lemma}{}{thm:rat-func-deg}
  Let $K=F(X)$ be the field of rational functions in $X$ over $F$. 
  If $u\in K$ with $u\notin F$ let $u=\nicefrac{f}{g}$ where $f,g\in F[X]$ and $\gcd(f,g)=1$. 
  Let $L=F(u)$. 
  Then $$[K\colon L]=\max\{\deg(f),\deg(g)\}$$
\end{lemma}
\begin{proof}
  By \SimpleExtthm we need to find the minimal polynomial of $X$ over $L$ to determine $[K\colon L]$. 
  Consider the polynomial $p(Y)=u\cdot g(Y)-f(Y)\in L[Y]$. 
  Then $X$ is a root of $p$. 
  Therefore, $X$ is algebraic over $L$ and so $[K\colon L]<\infty$ and $K=L(X)$. 

  Notice that $\deg(p)=\max\{\deg(f),\deg(g)\}$. 
  Let $f=\sum_{i=0}^n a_iX^i$ and $g=\sum_{j=0}^m b_i\cdot X^i$. 
  We will show that $p$ is irreducible over $L$. 
  Now the element $u$ cannot be algebraic over $k$ otherwise by \TowerLawthm we have $[K\colon F]=[K\colon L]\cdot [L\colon F]<\infty$, and we know $F(X)$ is an infinite extension of $F$. 
  So $u$ is transcendental over $F$ and hence $F[u]\cong F[X]$. 
  Now observe that the polynomial $p\in F[Y][u]\subseteq F(Y)[u]$ which is of degree $1$. 
  Hence, $p$ is irreducible over $k(Y)$. 
  Now $p\in k[Y][u]=k[u][Y]$ and therefore $p$ is irreducible over $k(u)$. 
  Hence, $p$ is the minimal polynomial of $X$ over $F$. 
  So we have the lemma.
\end{proof}

\begin{note}
In the proof of \TowerLawthm the proof do not assume that the field extensions needs to be finite. 
\end{note}

We will now prove a converse theorem of \FiniteAlglem. 

\begin{Theorem}{}{thm:fin-gen-alg}
  Let $K/F$ be a field extension. 
  If $\alpha_1,\ldots,\alpha_n\in K$ are algebraic over $F$, then
  \[
    [F(\alpha_1,\ldots,\alpha_n)\colon F] \;\leq\; \prod_{i=1}^{n}[F(\alpha_i)\colon F].
  \]
\end{Theorem}

\begin{proof}
  We will prove this by induction on $n$. 
  The base case, $n=1$ follows from the \TowerLawthm. 
  So suppose this is true for $n-1$. 
  Take $L=F(\alpha_1,\ldots,\alpha_{n-1})$. 
  Then by induction hypothesis $[L:F]\leq\prod_{i=1}^{n-1}[F(\alpha_i):F]$. 
  Since $\alpha_n$ is algebraic over $F$ and hence over $L$, with $\min(L,\alpha_n)\mid\min(F,\alpha_n)$ in $L[X]$, we have $[L(\alpha_n):L]\leq[F(\alpha_n):F]$. 
  Therefore, $[F(\alpha_1,\dots, \alpha_n)\colon L]\leq [F(\alpha_n)\colon F]$. 
  So by \TowerLawthm we have 
  \[
    [F(\alpha_1,\ldots,\alpha_n):F] = [L(\alpha_n):L]\cdot[L:F] \leq [F(\alpha_n):F]\cdot\prod_{i=1}^{n-1}[F(\alpha_i):F].
  \]
  So we have the theorem. 
\end{proof}

The inequality above can be strict in some cases. 
For example let $a=\sqrt[4]{2}$ and $b=\sqrt[4]{18}$. 
Then both $[\bbQ(a)\colon \bbQ]=[\bbQ(b)\colon \bbQ]=4$ as $X^4-2$ and $X^4-18$ are both irreducible polynomials over $\bbQ$. 
But $[\bbQ(\sqrt[4]{2},\sqrt3)\colon \bbQ]=8$ and $\sqrt[4]{18}\in  \bbQ(\sqrt[4]{2},\sqrt3)$. 
Hence, in this the inequality doesn't follow. 

So from \SimpleExtthm and \FinGenAlgthm we have the following criterion for an element being algebraic over a field. 

\begin{corolary}{}{cor:alg-degree}
  Let $K/F$ be a field extension and $\alpha\in K$. 
  Then $\alpha$ is algebraic over $F$ if and only if $[F(\alpha):F]<\infty$. 
  Moreover, $K$ is algebraic over $F$ if $[K:F]<\infty$.
\end{corolary}

Now we can extend \FinGenAlgthm to the case of fields generated by arbitrary number of elements or extending field.
\begin{Theorem}{}{thm:alg-gen}
  Let $K$ be a field extension of $F$ and $X\subseteq K$ such that each element of $X$ is algebraic over $F$. 
  Then $F(X)/F$ is algebraic. 
  If $|X|<\infty$, then $[F(X):F]<\infty$.
\end{Theorem}

\begin{proof}
  For any $\alpha\in F(X)$, \thmref{thm:fin-gen-sub} gives $\alpha\in F(a_1,\ldots,a_m)$ for some $a_1,\ldots,a_m\in X$. By \thmref{thm:fin-gen-alg}, $[F(a_1,\ldots,a_m):F]<\infty$, so $\alpha$ is algebraic over $F$ by \thmref{thm:finite-alg}. If $X$ is finite, the same bound gives $[F(X):F]<\infty$.
\end{proof}

\begin{Theorem}{Transitivity of Algebraic Extensions}{thm:alg-transitive}
  Let $F\subseteq L\subseteq K$ be fields. If $L/F$ and $K/L$ are algebraic, then $K/F$ is algebraic.
\end{Theorem}
\begin{proof}
  Let $\alpha\in K$. Since $K/L$ is algebraic, there is a nonzero $p(x)=\sum_{k=0}^m a_kx^k\in L[x]$ with $p(\alpha)=0$. Each $a_k$ is algebraic over $F$, so by \thmref{thm:fin-gen-alg} the extension $L'=F(a_0,\ldots,a_m)$ satisfies $[L':F]<\infty$. Since $p(x)\in L'[x]$ and $p(\alpha)=0$, the element $\alpha$ is algebraic over $L'$ and $[L'(\alpha):L']<\infty$. By the tower law,
  \[
    [F(a_0,\ldots,a_m,\alpha):F] = [L'(\alpha):L']\cdot[L':F] < \infty,
  \]
  so $\alpha$ is algebraic over $F$ by \thmref{thm:finite-alg}.
\end{proof}

%------------------------------------------------------------------
\subsection{Algebraic Closure and Transitivity}
%------------------------------------------------------------------

\begin{Definition}{Algebraic Closure in $K$}{alg-closure-def}
  Let $K$ be a field extension of $F$. The set
  \[
    \{\,a\in K : a\text{ is algebraic over }F\,\}
  \]
  is called the \emph{algebraic closure of $F$ in $K$}.
\end{Definition}

\begin{corolary}{}{cor:alg-closure-field}
  Let $K$ be a field extension of $F$ and let $L$ be the algebraic closure of $F$ in $K$. Then $L$ is a field, and therefore is the largest algebraic extension of $F$ contained in $K$.
\end{corolary}

\begin{proof}
  Let $a,b\in L$; both are algebraic over $F$. By \thmref{thm:fin-gen-alg}, $F(a,b)/F$ is finite, hence algebraic by \thmref{thm:finite-alg}. Since $a\pm b,\,ab,\,a/b\;(b\neq 0)$ all lie in $F(a,b)\subseteq L$, the set $L$ is closed under field operations, so it is a field. It is algebraic over $F$ by definition, and it contains any algebraic extension of $F$ inside $K$.
\end{proof}
